\section{Iteration names, projections and support proofs}
\label{annex:p06-forcing-iterations-support}

The recursive forcing inputs are
\cref{thm:p06-ftr-recursion,thm:p06-ftr-density,thm:p06-ftr-truth,thm:p06-ftr-zfc}.
All name and antichain choices are ground sets.

\begin{annexitem}{annex:p06.iter.name-bound}{thm:p06-iter-name-bound}{Mix only a bounded set of possible member names}
\begin{proof}
For a maximal antichain \(A\subseteq P\) and names \(\tau_a\), define
\[
 \operatorname{Mix}(A,\tau)=
 \{(\eta,q):(\exists a\in A)(\exists r)
       ((\eta,r)\in\tau_a\ \land q\leq a,r)\}.
\]
This is a name. Below \(a\), its supported pairs agree densely with
those of \(\tau_a\): the other antichain supports are incompatible
with \(a\), and a support \(r\) can be refined simultaneously with
\(a\) whenever it is compatible. Induction on the ranks of member
names, using the atomic membership/equality recursion, therefore gives
\(a\forces\operatorname{Mix}(A,\tau)=\tau_a\). In a supplied generic
containing \(a\), the same equality is also seen directly: a support
from another antichain member cannot lie in the filter, and common
strengthenings of \(a,r\) supply exactly the members of \(\tau_a^G\).
More explicitly for the atomic inclusion proof: a mixture pair supported
below \(b\in A\) can be active below \(a\) only when \(b=a\).
Its original support in \(\tau_a\) is then refined by every condition
being tested, so forced reflexivity of the member name witnesses its
membership in \(\tau_a\). Conversely every \((\eta,r)\in\tau_a\)
and every strengthening below \(a,r\) itself supplies a supported
\(\eta\)-pair in the mixture. These two dense membership checks are
exactly the atomic subset clauses, giving the claimed equality without
inferring forcing from an unqualified semantic quantification.

Let \(S\) be the set of member names occurring in the underlying set
name \(\dot Q\), together with \(\dot1_Q\). Form all mixtures over
ground maximal antichains of assignments into \(S\) for which
\(a\forces\tau_a\in\dot Q\) at each antichain member. These choices
are a set: antichains are subsets of \(P\), assignments are functions
into a fixed set, and each mixture is a subset of a fixed set of
member-name/support pairs. Include \(\dot1_Q\) itself. Density closure
shows that the unit forces each such mixture to belong to \(\dot Q\).
Their ranks are bounded by the ranks of names in \(S\).

If \(p\forces\tau\in\dot Q\), atomic membership gives densely below
\(p\) a condition forcing \(\tau\) equal to a member name from
\(\dot Q\), and forcing that member to belong to \(\dot Q\).
Choose a maximal antichain of these conditions below \(p\); complete
it to a global maximal antichain using conditions incompatible with
\(p\). Assign the selected member names on the first part and
\(\dot1_Q\) on the second. The mixture is in the bounded set, and
its equality with \(\tau\) is dense below \(p\), hence forced by
\(p\). This proves the representation claim without a proper class
of unrestricted condition names. The same argument supplies bounded
names for existential witnesses forced to belong to \(\dot Q\).
For that last assertion, choose an antichain of existential-witness
refinements, then refine each witness's membership in \(\dot Q\) to
equality with a name in \(S\). Mixing these \(S\)-names gives a
single bounded witness: on the antichain it has the required property,
and density closure forces that property below the starting condition.
For general existential names, Collection over the set antichain first
bounds the witness-name family. Choice selects that set family for the
mixing constructor. In particular the enumeration names used below do
not depend on an unstated choice from a proper class.
\end{proof}
\end{annexitem}

\begin{annexitem}{annex:p06.iter.preorder}{thm:p06-iter-preorder}{Quotient and the exact projection lifting property}
\begin{proof}
Reflexivity follows from forced reflexivity of the iterand order. If
\((s,\rho)\leq(r,\nu)\leq(p,\mu)\), then \(s\leq p\), and
monotonicity transfers \(r\forces\nu\leq\mu\) to \(s\); forced
transitivity gives \(s\forces\rho\leq\mu\). Thus the relation is
a preorder. Mutual extension is an equivalence relation; order on its
classes is well defined by transitivity and is antisymmetric.
Mutually extending pairs have mutually extending first coordinates,
which are equal because \(P\) is a partial order. Thus \(\pi\) is
well defined. A first-coordinate strengthening \(r\leq p\) lifts
to \((r,\mu)\leq(p,\mu)\), giving the asserted projection property.

The unit-coordinate map preserves and reflects order. A common lower
bound of images projects to a common lower bound in \(P\); conversely
the image of such a lower bound is common. A maximal antichain in
\(P\) is predense after the map: below any pair, choose a compatible
antichain member in its first coordinate and lift the common
strengthening with the unchanged second name.

In a general preorder, quotient compatibility is equivalent to original
compatibility by choosing a representative of a common class lower
bound. Dense sets map to dense class sets and their preimages are dense;
filters lift to unions of classes and project in the other direction.
Translate a name by replacing each support condition by its class and
recursing on member names. Conversely choose representatives of supports
and recurse. Rank induction gives the same generic evaluations; atomic
and formula induction gives the same forcing clauses because common
strengthenings and density are unchanged. This proves the quotient
bridge used at every subsequent iteration stage.
\end{proof}
\end{annexitem}

\begin{annexitem}{annex:p06.iter.two-generics}{thm:p06-iter-two-generics}{Actual filters, dense-set names and flattened names}
\begin{proof}
Use the lifted filter on the preorder. From a two-step generic \(H\),
put \(G=\{p:(p,\dot1_Q)\in H\}\). Projection and its lifting property
show that it is a filter and meets every ground dense set: the pairs
whose first coordinates belong to such a set are dense. Define
\(K=\{\mu^G:(p,\mu)\in H\}\). It is directed because \(H\) has
common stronger pairs and the forced inequalities hold in \(M[G]\).
To check weaker closure, represent an actual weaker condition by a
bounded total name \(\nu\), and choose a member of \(G\) forcing
the inequality, by truth. The corresponding unit-coordinate pair is
in \(H\); a common strengthening with the original pair is weaker
than a pair with right coordinate \(\nu\), so that pair belongs to
\(H\). Thus the weaker evaluated condition is in \(K\).
Every actual \(Q\)-member has such a bounded name: its occurrence in
the evaluation of the set name gives a ground member name and a generic
support, to which the representation lemma applies.

Let \(E\in M[G]\) be dense open in \(Q=\dot Q^G\), with ground
name \(\dot E\). Truth supplies \(p_0\in G\) forcing that property.
The following is a ground dense set of two-step conditions: pairs with
first coordinate incompatible with \(p_0\), and pairs \((p,\mu)\)
with \(p\leq p_0\) forcing \(\mu\in\dot E\). Indeed, a first
coordinate compatible with \(p_0\) can be strengthened below it;
there density of \(\dot E\), the existential witness clause and the
bounded-name lemma provide a stronger right coordinate in \(\dot E\).
Since \((p_0,\dot1_Q)\in H\), the incompatible option cannot meet
\(H\). Thus \(K\) meets \(E\), proving second-stage genericity.

The equality with the stated product filter requires its reverse inclusion.
If \(p\in G\), \(\mu^G\in K\), choose \((r,\nu)\in H\) with
\(\nu^G=\mu^G\). A member \(s\in G\) forces that equality; taking
a common strengthening in \(H\) of \((r,\nu)\), \((p,\dot1_Q)\)
and \((s,\dot1_Q)\) gives a condition below \((p,\mu)\). Weaker
closure therefore puts \((p,\mu)\) in \(H\).

Conversely, given the two generics, the displayed product filter is
directed. A common \(K\)-strengthening of two right values has a bounded
ground total name; truth and a common member of \(G\) force the needed
two inequalities. Weaker closure is immediate from the forced orders.
For ground dense open \(D\subseteq R\), define in \(M[G]\)
\(D^G=\{\mu^G:(\exists p\in G)(p,\mu)\in D\}\).
For every bounded total name \(\nu\), the ground set of \(p\) admitting
a \(D\)-pair \((p,\mu)\) with \(p\forces\mu\leq\nu\) is dense
in \(P\), by applying density below \((p,\nu)\). It meets \(G\).
Since every \(Q\)-condition is represented, \(D^G\) is dense in
\(Q\). The second generic meets it, giving a member of the product
filter in \(D\). Quotient and preorder genericity correspond by the
preceding bridge.

For completeness define the name translation. Let
\(\operatorname{pair}_P(u,v)\) be the usual \(P\)-name for the ordered
pair of evaluations, constructed from the two singleton/Pairing names.
By rank recursion set
\[
 \tau^\downarrow=
 \{(\operatorname{pair}_P(\sigma^\downarrow,\mu),p):
                         (\sigma,(p,\mu))\in\tau\}.
\]
After evaluation by \(G\), this is a \(Q\)-name: its condition
coordinates are \(\mu^G\in Q\), and its member names have smaller
rank. Induction gives
\((\tau^\downarrow{}^G)^K=\tau^H\), because a support contributes
exactly when \(p\in G\) and \(\mu^G\in K\). Quotient names first
lift recursively to preorder supports. Conversely \(G,K\) are definable
from \(H\), and \(H\) is a set definable from \(G,K\). The transitive
generated-extension/name-evaluation property therefore gives both
inclusions in \(M[H]=M[G][K]\), with the translation providing the
explicit forward name construction.
\end{proof}
\end{annexitem}

\begin{annexitem}{annex:p06.iter.initial-segments}{thm:p06-iter-initial-segments}{Recursive sets, quotients and tail amalgamation}
\begin{proof}
At each stage the bounded name lemma supplies a set of coordinate names.
At a limit the raw functions form a subset of the product of these sets;
the finite or countable support test and initial-segment tests are
Separation conditions. Transfinite recursion and Replacement give the
set system through the set ordinal \(\delta\).
Reflexivity and transitivity of each raw order follow by induction:
an earlier restriction of a stronger condition is stronger at every
earlier stage, so monotonicity transfers the coordinate inequalities,
and iterand transitivity composes them. The quotient is a partial order.

Mutual raw extension gives mutual extension of every restriction; thus
the restriction classes do not depend on representatives. Extending by
units preserves mutual extension and is also well defined. Composition
of restrictions and unit extensions is literal composition of functions.
For the lifting property choose a raw representative \(s\) of \(r\)
and a representative \(p\) of the longer condition. Replace the prefix
of \(p\) by \(s\). Induction on the tail coordinates proves that
the resulting prefix is stronger than the old prefix: at each new
coordinate the coefficient is unchanged and its reflexive inequality
is forced. Its support is contained in the union of two allowed
supports, so the result is a condition with exactly the restriction \(r\).

The unit embedding preserves and reflects order and incompatibility.
To see completeness, take a maximal antichain at the shorter stage;
below any longer condition choose a compatible short antichain member,
then lift the common short strengthening by the preceding construction.
The images are predense. Their restrictions therefore give generic
filters on initial stages. Two-step factoring supplies successive
generic extensions; no claim about a fresh support convention for an
arbitrary tail quotient is needed for this conclusion.

Finally suppose \(p,q\) have no common nonunit support coordinate
above \(\alpha\), and their restrictions have a common strengthening
\(r\). Keep \(r\) below \(\alpha\), and above it use the coefficient
of \(p\) or \(q\) at each of their disjoint supports, with units
elsewhere. At a coordinate from \(p\), its preceding combined restriction
already strengthens the prefix of \(p\), and the unchanged coefficient
is therefore valid and stronger; the coefficient of \(q\) there is
the greatest condition. The same argument applies at a \(q\)-coordinate.
Induction proves both inequalities. The union support remains allowed.
This is the asserted disjoint-tail amalgamation, not a raw coordinatewise
union over shared incompatible names.
\end{proof}
\end{annexitem}

\begin{annexitem}{annex:p06.iter.two-ccc}{thm:p06-iter-two-ccc}{Bounding the indices of a possible antichain}
\begin{proof}
First a ccc forcing covers a name for a countable set of ground ordinals
by a ground countable set. Choose a name for an enumeration, using the
existential witness/mixing construction. At each natural coordinate
choose a maximal antichain deciding its ordinal value; it is countable.
The union of all possible values is a ground countable set, and density
forces that it covers the enumerated set. If the original assertion is
forced only below \(p\), perform the argument below \(p\).
The same possible-value argument shows that ccc preserves \(\omega_1\).

If \(\{(p_\xi,\mu_\xi):\xi<\omega_1\}\) were a two-step antichain,
let \(\dot I\) name the indices for which \(p_\xi\) enters the first
generic. For \(\xi,\eta\in I\), the evaluated right coordinates
must be incompatible: a compatible right strengthening, named and forced
below a common member of the first generic, would give a two-step
common strengthening. They must also be distinct. Thus \(P\) forces
the indexed right family to be an antichain in \(\dot Q\), hence
\(\dot I\) countable. The preceding claim can be checked inside the
forcing calculation: a condition forcing a compatible pair of indexed
entries can be refined to two distinct check indices, and then below
both of their canonical generic supports. The bounded witness lemma
names a common right condition, giving a ground two-step common
strengthening, impossible for the given antichain. Equality of two
indexed values gives the same contradiction. Decision density therefore
forces the antichain and injectivity claims before forced ccc is applied;
no vacuous universal quantification over nonexistent generics is used.
The covering argument gives a ground countable
\(B\subseteq\omega_1\) with \(1_P\forces\dot I\subseteq\check B\).
But for \(\xi\notin B\), \(p_\xi\) forces \(\xi\in\dot I\)
by the canonical generic name, a contradiction. Quotienting does not
change the antichain conclusion.
\end{proof}
\end{annexitem}

\begin{annexitem}{annex:p06.iter.fs-ccc}{thm:p06-iter-fs-ccc}{Finite delta systems and tail amalgamation}
\begin{proof}
We supply the finite-set lemma. An uncountable indexed family of finite
sets has an uncountable subfamily with a fixed intersection root. First
restrict to one finite size and induct on that size. If some element
belongs to uncountably many sets, restrict to them, remove it, apply
the induction hypothesis and add it to the root. If each element belongs
to only countably many indexed sets, recursively for \(\xi<\omega_1\)
choose the next set disjoint from all previously chosen sets. Their
union at each stage is countable, and only countably many indexed sets
meet it, so an available set remains. This gives an uncountable disjoint
family, with empty root. The empty-size case is immediate, including
repeated supports.

Induct on iteration length. The unit stage is ccc; successors use the
proved two-step theorem. At a limit \(\delta\), choose raw representatives
of a proposed uncountable antichain and apply the lemma to their finite
supports, obtaining a fixed finite root. Choose \(\alpha<\delta\)
above that root. Among the uncountably many restrictions to \(P_\alpha\),
either two are equal, or there are uncountably many distinct ones and
the induction hypothesis gives two compatible restrictions. The tails
of that pair have disjoint supports above \(\alpha\), since their
intersection was precisely the root. The proved tail-amalgamation
construction makes the full pair compatible, contradicting the
antichain. This completes every ordinal-length case.
\end{proof}
\end{annexitem}

\begin{annexitem}{annex:p06.iter.cs-closed}{thm:p06-iter-cs-closed}{Names for coordinate lower bounds on one countable support}
\begin{proof}
Let \(p_{n+1}\leq p_n\) be a descending sequence of raw representatives;
if the sequence is given as quotient conditions, choose representatives,
whose order comparisons still hold. Let \(S\) be the countable union
of their supports. Construct a raw \(r\) on \(S\) in increasing
coordinate order. Suppose its earlier restrictions strengthen every
earlier restriction of each \(p_n\). They force the coordinate sequence
\(\langle p_n(\xi):n<\omega\rangle\) to be descending in
\(\dot Q_\xi\), by monotonicity of all the earlier comparisons.
The stage hypothesis forces this sequence to have a lower bound.
The existential witness and bounded mixed-name lemma therefore supply
\(r(\xi)\in\mathcal B_\xi\) forced by that earlier restriction to
be below every \(p_n(\xi)\). At a coordinate outside \(S\), use the
unit. At limit coordinates, the raw earlier function has countable
support and belongs to the already defined limit stage; no lower bound
of arbitrary proper conditions is assumed.

The resulting \(r\) is a condition: its support is included in \(S\),
all coefficient names are total valid members, and induction gives all
required earlier-stage comparisons. It strengthens every \(p_n\).
Thus the countable-support iteration is countably closed, and the
current complete countable-closure master construction
\cref{thm:p06-pfm-closed} makes it proper. This proof applies to
countably closed iterands; an arbitrary proper iterand need not provide
the coordinate lower bound used here.
\end{proof}
\end{annexitem}

\begin{annexitem}{annex:p06.iter.proper-cover}{thm:p06-iter-proper-cover}{Countable possible values below a master condition}
\begin{proof}
Suppose \(p\) forces a name \(\dot f\) to map \(\omega\) into a
ground ordinal \(\lambda\). For each natural coordinate choose a
name for its value, and a maximal antichain deciding that value, with
the corresponding ground ordinal-value map. Choose a sufficiently
large regular \(\chi\), and a countable \(N\prec H_\chi\) containing
\(p,P,\lambda,\dot f\) and the countable sequences of these antichains
and value maps. Every individual antichain and map belongs to \(N\),
by the countable-element closure already proved in
\cref{thm:p06-pfm-hulls}. Conditional value names can first be mixed
with a default value outside \(p\), and antichains made global, so all
tests are actual maximal-antichain tests.

Properness gives \(q\leq p\) master for \(N\). For each coordinate
its antichain's intersection with \(N\) is predense below \(q\).
A member in \(N\) has its decided value in \(N\), since its value
map and the member both belong to the elementary model. Hence the
conditions forcing that coordinate into \(N\cap\lambda\) are dense
below \(q\), and density forces that membership. All evaluated values
therefore lie in the ground countable set \(N\cap\lambda\).

Perform this construction below every condition forcing the function
assertion; such countable-cover conditions are dense below \(p\).
A generic containing \(p\) meets them. A name for a countable set of
ordinals has a name for an enumeration and a ground ordinal bound on
its values, by the ordinary name-rank bound. Thus every such set in the
extension has a ground countable cover. The empty set requires none
of this construction. For \(\lambda=\omega_1\), the ground cover is
bounded below \(\omega_1\), so no new countable cofinal enumeration
can collapse \(\omega_1\). This proves preservation.

A Prikry generic sequence is a new countable unbounded subset of the
ground regular \(\kappa\). Any ground countable cover inside
\(\kappa\) would be bounded there. The preceding conclusion therefore
contradicts the actual Prikry cofinal sequence, proving nonproperness
for both preceding Prikry presentations despite their cardinal preservation.
\end{proof}
\end{annexitem}

\begin{annexitem}{annex:p06.iter.support-examples}{thm:p06-iter-support-examples}{Atomic stages make the support distinction explicit}
\begin{proof}
Let \(A_\theta\) have a greatest condition and \(\theta\) mutually
incompatible atoms below it. A descending sequence either stays at the
unit, which has an atomic lower bound, or reaches an atom and remains
there. Hence it is countably closed and proper. Its generic extension
adds no sets: the dense atom set is met by some ground atom, and the
resulting principal filter belongs to the ground.

For a finite iteration of constant \(A_{\theta_n}\), decide the names
successively below a complete earlier atomic choice, and choose an atom
at the next coordinate. This gives a dense set of complete finite
atomic sequences, with the evident extension order. In a finite-support
length-\(\omega\) iteration, any condition has bounded finite support;
the same finite procedure gives a stronger complete finite initial
sequence. Thus its dense presentation is the finite initial sequences
\(s\), with \(s(n)<\theta_n\). If all \(\theta_n=2\), this is
precisely Cohen \(2^{<\omega}\).

In countable support, recursively choose atomic values at every natural
coordinate, below the name evaluated by the earlier finite atomic
prefix. The resulting entire ground sequence is an allowed condition
stronger than the original one. It is a minimal quotient condition:
any strengthening must have the same atomic value at every coordinate,
and is equivalent. Entire atomic sequences are therefore dense.
For \(\theta_n=2\), the \(2^{\aleph_0}\) different entire sequences
are pairwise incompatible, so this countable-support forcing is not ccc,
although it is atomic and proper and adds no sets.

Now put \(\theta_n=\aleph_{n+1}\) and
\(\kappa=\sup_n\theta_n=\aleph_\omega\), all ground values.
In the finite-support dense presentation, the generic union is a function
\(f\) with domain \(\omega\), \(f(n)<\theta_n\). For each
\(\alpha<\kappa\), requiring \(f(n)=\alpha\) at some coordinate
is dense: choose an unused sufficiently late coordinate with
\(\theta_n>\alpha\), then extend the finite sequence to it.
Genericity consequently makes \(f\) onto the entire ground \(\kappa\).
It becomes countable; in particular ground \(\omega_1\) collapses,
so the limit is not proper. Every stage was nevertheless countably
closed, in fact atomic, so was forced proper at every finite stage.
The countable-support version again has dense entire ground atomic
sequences and changes no sets or cardinals. This is an actual infinite
support failure, not a finite example promoted to the full proper theorem.
\end{proof}
\end{annexitem}

\begin{annexitem}{annex:p06.iter.native-bridge}{thm:p06-iter-native-bridge}{Retain the Magidor clauses under order-dual transport}
\begin{proof}
The sets of strengthening conditions coincide by the explicit definition
\(r\leq_g p\) iff \(r\geq p\), and likewise for direct extension.
Thus common strengthenings, compatibility and antichains coincide.
Native greater-directed/weakly closed filters become host lower-directed/
upward-closed filters, meeting the identical dense sets. The weakest
native condition is the greatest host condition. The condition pairs
in name constructors do not change; induction on name rank gives
identical evaluations. Atomic recursive forcing quantifies over the
same strengthening and density domains, so name-rank induction gives
the same atomic values; formula induction then gives the same negation
and quantifier values.

Throughout this identity transport, the native full coordinate support
and the finite set of coordinates permitting nondirect changes remain
in the order predicate. Erasing either would change the strengthening
set and invalidate the proof. Thus the orientation bridge supplies
no claim of countable-support proper preservation for this different
native Prikry-type iteration.
\end{proof}
\end{annexitem}
